% 高级 QED：真空极化势、软辐射与强平面波中的精确 Dirac 解

同一 QED 母理论在三个不同极限中呈现出三种可复用结构：真空极化修正静态光子传播子，软定理控制低能真实光子辐射，经典平面波背景则可以非微扰地吸收到 Dirac 方程的零阶解中。三者都沿用物理四动量、$q_{\mathrm e}=-e$ 和统一 Fourier 约定。

前面的圈图计算说明，虚粒子会修正传播子本身。要看这种修正在可观测空间势中留下什么痕迹，最直接的做法是先取静态极限：此时真空极化只改变 Coulomb 光子核的动量依赖，再通过三维 Fourier 变换回到位置空间。静态 Coulomb 核的一圈真空极化修正可写成
\begin{equation}
 \delta\widetilde V_U(\mathbf q)
 =\widetilde V_C(\mathbf q)\widehat\Pi(Q^2),
 \label{eq:Uehling-momentum-correction-r12}
\end{equation}
其中类空动量满足 $Q^2\equiv\mathbf q^2>0$。重整化真空极化的谱表示为
\begin{equation}
 \frac{\widehat\Pi(Q^2)}{Q^2}
 =\frac{\alpha}{3\pi}
 \int_{4m_{\mathrm e}^2c^2}^{\infty}\frac{\dd s}{s(s+Q^2)}
 \sqrt{1-\frac{4m_{\mathrm e}^2c^2}{s}}
 \left(1+\frac{2m_{\mathrm e}^2c^2}{s}\right).
 \label{eq:Pi-dispersion-explicit-r12}
\end{equation}
对三维动量作 Fourier 变换后得到
\begin{equation}
 \boxed{
 \delta V_U(r)
 =V_C(r)\frac{2\alpha}{3\pi}
 \int_1^\infty\dd t\,
 \ee^{-2rt/\lambda_C}
 \left(1+\frac1{2t^2}\right)
 \frac{\sqrt{t^2-1}}{t^2},
 \qquad
 \lambda_C\equiv\frac{\hbar}{m_{\mathrm e} c}.}
 \label{eq:Uehling-potential-SI-r12}
\end{equation}
积分核为正；对电子与正核的吸引 Coulomb 势 $V_C<0$，因此 $\delta V_U<0$。修正只在 $r\lesssim\lambda_C$ 的短距离显著，并以一圈参数 $\alpha/\pi$ 控制；更高圈从相对 $O(\alpha/\pi)$ 继续修正该一圈结果。

\begin{figure}[htbp]
 \centering
 \begin{tikzpicture}[scale=.82,baseline=(current bounding box.center),
 fermion/.style={thick,postaction={decorate},decoration={markings,mark=at position .58 with {\arrow{Latex}}}},
 photon/.style={thick,decorate,decoration={snake,amplitude=1.25pt,segment length=6pt}}]
\node[draw,circle,inner sep=2pt] (N) at (-3,0) {$Ze$};
\node[draw,circle,inner sep=2pt] (e) at (3,0) {$e^-$};
\draw[photon] (N)--(-1.0,0);
\draw[photon] (1.0,0)--(e);
\draw[fermion] (-1.0,0) arc (180:0:1.0 and .85);
\draw[fermion] (1.0,0) arc (0:-180:1.0 and .85);
\node[above] at (0,1.15){vacuum polarization};
\node[below] at (0,-1.15){$q^0=0$};
\end{tikzpicture}

 \caption{静态 Coulomb 交换中插入一次电子--正电子真空极化圈。重整化自能改变光子传播子的谱权重，三维 Fourier 变换后得到\eqnrefs{eq:Uehling-potential-SI-r12} 的 Uehling 势。}
 \label{fig:uehling-loop-r12}
\end{figure}

\begin{figure}[htbp]
 \centering
 \includegraphics[width=.73\linewidth]{uehling_relative.png}
 \caption{点核 Uehling 修正的相对大小 $\delta V_U/V_C$ 随 $r/\lambda_C$ 的变化。曲线由\eqnrefs{eq:Uehling-potential-SI-r12} 数值积分得到。}
 \label{fig:uehling-relative-r12}
\end{figure}

真空极化处理的是虚光子圈对传播的修正；真实光子辐射则引入另一种尺度分离。当额外辐射光子的能量远小于硬散射的特征能量时，完整含光子振幅的领先项不再依赖硬子图的细节，而只由外部带电粒子的运动学决定。设电子在一个硬散射过程中从速度 $\boldsymbol\beta_{{\rm rel},i}$ 变到 $\boldsymbol\beta_{{\rm rel},f}$，并额外辐射方向为 $\mathbf n$、角频率为 $\omega$ 的软光子。若 $E_h$ 表示硬过程的最小特征能量，则
\begin{equation*}
 \epsilon_{\rm soft}^{(\gamma)}\equiv\frac{\hbar\omega}{E_h}\ll1
\end{equation*}
控制软展开。保留 $1/\omega$ 的领先外腿因子后，单光子辐射概率为
\begin{equation}
 \boxed{
 \dd P_{\rm soft}
 =\frac{\alpha}{4\pi^2}
 \frac{\dd\omega}{\omega}\,\dd\Omega
 \left|
 \frac{\boldsymbol\beta_{{\rm rel},f}^\perp}{1-\mathbf n\cdot\boldsymbol\beta_{{\rm rel},f}}
 -\frac{\boldsymbol\beta_{{\rm rel},i}^\perp}{1-\mathbf n\cdot\boldsymbol\beta_{{\rm rel},i}}
 \right|^2.}
 \label{eq:soft-photon-probability-r12}
\end{equation}
相对于领先软项，有限光子能量修正由 $\epsilon_{\rm soft}^{(\gamma)}$ 控制；当它不再远小于一时，必须恢复完整含光子过程的矩阵元。多软光子与虚软修正的包容处理由一般红外母式给出，不能把\eqnrefs{eq:soft-photon-probability-r12} 单独积分到 $\omega=0$。

\begin{figure}[htbp]
 \centering
 \begin{tikzpicture}[scale=.78, baseline=(current bounding box.center), hard/.style={draw,rounded corners,minimum width=1.0cm,minimum height=.9cm,fill=black!4}]
\begin{scope}[xshift=-4cm]
\node[hard] (H) at (0,0) {$\mathcal H$};
\draw[fermion] (-2.2,0)--(-.6,0);\draw[fermion] (.6,0)--(2.2,0);
\coordinate (v) at (1.25,0);
\draw[photon] (v)--(2.1,1.35);
\node[below] at (1.45,-.4){outgoing};
\end{scope}
\node at (0,0){$+$};
\begin{scope}[xshift=4cm]
\node[hard] (H) at (0,0) {$\mathcal H$};
\draw[fermion] (-2.2,0)--(-.6,0);\draw[fermion] (.6,0)--(2.2,0);
\coordinate (v) at (-1.25,0);
\draw[photon] (v)--(-2.1,1.35);
\node[below] at (-1.45,-.4){incoming};
\end{scope}
\end{tikzpicture}

 \caption{软光子附着在同一条外部电子线的出射与入射两种位置。领先 $1/\omega$ 奇异性只依赖外腿四动量，因此从任意硬子图因子化。}
 \label{fig:soft-brems-r12}
\end{figure}

\begin{figure}[htbp]
 \centering
 \includegraphics[width=.76\linewidth]{soft_brems_angular.png}
 \caption{软光子辐射的角向几何因子。取入射与出射电子速率均为 $\beta_{\rm rel}=0.8$，电子偏转角为 $25^\circ$；曲线只计算\eqnrefs{eq:soft-photon-probability-r12} 方括号中的几何因子，因此不对应特定实验装置。}
 \label{fig:soft-brems-angular-r12}
\end{figure}

软展开仍把电磁场作为量子辐射处理；若外加光场具有很高占据数，而且其经典波形已知，更自然的组织方式是把这部分场直接并入零阶 Dirac 方程，而不是把它逐光子展开。平面波背景是少数能够把这一步精确完成的时空依赖场。令经典四势只依赖无量纲相位
\begin{equation*}
 \phi\equiv\kappa\cdot x,
 \qquad
 \kappa^2=0,
 \qquad
 \kappa\cdot A(\phi)=0,
\end{equation*}
其中 $\kappa^\mu$ 是逆长度量纲的平面波四波矢。背景中的 Dirac 方程为
\begin{equation}
 [\ii\hbar\slashed\partial-q_{\mathrm e}\slashed A(\phi)-m_{\mathrm e} c]\psi(x)=0.
 \label{eq:Volkov-Dirac-r12}
\end{equation}
取自由在壳四动量 $p^\mu$ 和自旋子 $u_s(p)$，并定义
\begin{align*}
 R_p(\phi)
 &\equiv
 1+\frac{q_{\mathrm e}}{2p\cdot\kappa}
 \slashed\kappa\slashed A(\phi),\\
 F_p(\phi)
 &\equiv
 \frac{q_{\mathrm e} p\cdot A(\phi)}{p\cdot\kappa}
 -\frac{q_{\mathrm e}^2A^2(\phi)}{2p\cdot\kappa},\\
 S_p(x)
 &\equiv p\cdot x+\int^{\phi}\dd\varphi\,F_p(\varphi).
\end{align*}
$R_p$ 无量纲，而 $F_p$ 与 $S_p$ 具有作用量量纲，因此 $S_p/\hbar$ 无量纲。任意强度平面波中的精确 Volkov 解为
\begin{equation}
 \boxed{
 \psi_{p,s}^{\rm V}(x)
 =R_p(\phi)u_s(p)
 \exp\left[-\frac{\ii}{\hbar}S_p(x)\right].}
 \label{eq:Volkov-solution-r12}
\end{equation}
它不是弱场展开；背景强度全部保留在 $R_p$ 与 $S_p$ 中。

若背景对 $\phi$ 周期，并选取 $\langle A\rangle_\phi=0$ 的规范，周期平均相位定义准动量
\begin{equation}
 \boxed{
 \bar p^\mu
 =p^\mu
 -\frac{q_{\mathrm e}^2\langle A^2\rangle_\phi}{2p\cdot\kappa}\kappa^\mu,
 \qquad
 \bar p^2=m_{\mathrm e}^2c^2-q_{\mathrm e}^2\langle A^2\rangle_\phi.}
 \label{eq:volkov-quasimomentum-dp43}
\end{equation}
以峰值矢势定义无量纲强度
\begin{equation*}
 a_0\equiv\frac{|q_{\mathrm e}|A_0}{m_{\mathrm e} c}.
\end{equation*}
圆偏振与线偏振分别满足
\begin{equation}
 \boxed{
 m_*^{\rm circ}=m_{\mathrm e}\sqrt{1+a_0^2},
 \qquad
 m_*^{\rm lin}=m_{\mathrm e}\sqrt{1+\frac{a_0^2}{2}}.}
 \label{eq:volkov-effective-mass-dp43}
\end{equation}
只有进一步取 $a_0\ll1$ 时，才可把这些精确平面波结果重新展开成普通弱场多光子微扰级数。

\begin{figure}[htbp]
 \centering
 \begin{tikzpicture}[scale=.82,baseline=(current bounding box.center),
 fermion/.style={very thick,postaction={decorate},decoration={markings,mark=at position .6 with {\arrow{Latex}}}},
 wave/.style={thick,decorate,decoration={snake,amplitude=2.0pt,segment length=8pt}}]
\draw[fermion] (-3.2,-.8) .. controls (-1.2,-.2) and (1.2,.2) .. (3.2,.8);
\foreach \y in {-1.2,-.6,0,.6,1.2}{\draw[wave] (-2.8,\y+1.0)--(2.8,\y-1.0);}
% Labels are intentionally outside the dense wave bundle.
\node[anchor=east,fill=white,inner sep=1.2pt] at (-3.0,-.92){$p^\mu$};
\node[anchor=west,fill=white,inner sep=1.2pt] at (3.0,.92){$\bar p^\mu$};
\node[anchor=south,fill=white,inner sep=2pt] at (0,2.12){classical background $A^\mu(\kappa\!\cdot\! x)$};
\end{tikzpicture}

 \caption{Volkov 问题的物理结构：电子在任意强度经典平面波背景中传播。波线表示外加经典背景，不是内部光子传播子。}
 \label{fig:volkov-plane-wave-r12}
\end{figure}

\begin{figure}[htbp]
 \centering
 \includegraphics[width=.72\linewidth]{volkov_effective_mass.png}
 \caption{Volkov 有效质量随归一化矢势 $a_0$ 的变化。线偏振与圆偏振的差异来自周期平均 $\langle A^2\rangle_\phi$ 不同；曲线直接使用\eqnrefs{eq:volkov-effective-mass-dp43}。}
 \label{fig:volkov-effective-mass-r12}
\end{figure}

Uehling 势、软制动辐射与 Volkov 解分别对应 QED 的三个受控读取：一圈传播子谱修正、低能真实光子的外腿因子化，以及强经典平面波被精确吸收到零阶 Dirac 算符。它们共享同一作用量和 SI 约定，差异只在所选择的尺度层级与外场处理方式。

